\section{Παράρτημα Α: Πλήρης Κώδικας Υλοποίησης}
 
Το παρόν παράρτημα περιέχει τον πλήρη κώδικα Python που αναπτύχθηκε στο πλαίσιο της διατριβής. Ο κώδικας οργανώνεται σε πέντε Jupyter Notebooks, τα οποία παρουσιάζονται στις ακόλουθες πέντε υποενότητες.
 
\textbf{Περιβάλλον εκτέλεσης:}
\begin{itemize}
\item Python 3.12
\item \texttt{numpy}, \texttt{pandas}, \texttt{scipy}: Βασικές βιβλιοθήκες
\item \texttt{scikit-learn}: Tree-based models, evaluation
\item \texttt{pytorch}: Deep learning (RNN, LSTM)
\item \texttt{xgboost}, \texttt{lightgbm}: Gradient boosting
\item \texttt{imbalanced-learn}: SMOTE
\item \texttt{matplotlib}, \texttt{seaborn}: Visualization
\end{itemize}
 
Για πλήρη αναπαραγωγιμότητα, όλα τα πειράματα χρησιμοποιούν σταθερό \texttt{random\_state=42}.
 
 
 
